\documentclass[12pt,a4paper,oneside]{report}
\usepackage[utf8]{inputenc}
\usepackage[T1]{fontenc}
\usepackage[french]{babel}
\usepackage{amsmath,amssymb,amsfonts}
\usepackage{graphicx}
\usepackage{xcolor}
\usepackage{hyperref}
\usepackage{booktabs}
\usepackage{multirow}
\usepackage{geometry}
\usepackage{float}
\usepackage{tcolorbox}
\usepackage{enumitem}
\usepackage{listings}
\usepackage{caption}

\geometry{margin=2.5cm}

% Couleurs Institutionnelles & Modernes
\definecolor{navyblue}{RGB}{15, 23, 42}
\definecolor{royalblue}{RGB}{30, 58, 138}
\definecolor{accentblue}{RGB}{37, 99, 235}
\definecolor{lightbg}{RGB}{248, 250, 252}
\definecolor{emeraldgreen}{RGB}{16, 185, 129}
\definecolor{marsaorange}{RGB}{234, 88, 12}

\hypersetup{
    colorlinks=true,
    linkcolor=royalblue,
    citecolor=accentblue,
    urlcolor=accentblue
}

\newtcolorbox{bestbox}[1]{
    colback=emeraldgreen!10!white,
    colframe=emeraldgreen!80!black,
    fonttitle=\bfseries\color{white},
    title=#1,
    arc=4pt
}

\newtcolorbox{infobox}[1]{
    colback=royalblue!5!white,
    colframe=royalblue!80!black,
    fonttitle=\bfseries\color{white},
    title=#1,
    arc=4pt
}

\newtcolorbox{marsabox}[1]{
    colback=marsaorange!5!white,
    colframe=marsaorange!85!black,
    fonttitle=\bfseries\color{white},
    title=#1,
    arc=4pt
}

\lstset{
    basicstyle=\ttfamily\small,
    keywordstyle=\color{royalblue}\bfseries,
    stringstyle=\color{emeraldgreen},
    commentstyle=\color{gray}\itshape,
    breaklines=true,
    frame=single,
    backgroundcolor=\color{lightbg}
}

\begin{document}

% ==============================================================================
% PAGE DE GARDE
% ==============================================================================
\begin{titlepage}
    \centering
    
    % Logos Institutionnels (EMSI à gauche, Marsa Maroc à droite)
    \begin{minipage}{0.42\textwidth}
        \centering
        \includegraphics[width=0.85\textwidth,height=2.2cm,keepaspectratio]{logo_emsi.png}
    \end{minipage}
    \hfill
    \begin{minipage}{0.42\textwidth}
        \centering
        \includegraphics[width=0.85\textwidth,height=2.2cm,keepaspectratio]{logo_marsa.png}
    \end{minipage}

    \vfill

    % Titres et lignes de séparation
    \rule{\linewidth}{1.5pt}\\[0.3cm]
    {\LARGE \bfseries Exploitation et Logistique Portuaire :} \\[0.3cm]
    {\Large \bfseries Conception d'un Tableau de Bord Analytique et Intégration d'un Système RAG pour le Suivi des Flux à Marsa Maroc}\\[0.2cm]
    \rule{\linewidth}{1.5pt}

    \vfill

    {\large \textbf{RAPPORT DE STAGE D'IMMERSION PROFESSIONNELLE}}\\[0.15cm]
    \textit{Stage d'immersion professionnelle d'un mois et demi}

    \vfill

    % Auteurs et Encadrantes
    \begin{minipage}{0.48\textwidth}
        \begin{flushleft}
            \textbf{Réalisé par :}\\[0.15cm]
            \textbf{NASSIH Houssam}\\
            \textbf{NASSIH Anas}\\[0.1cm]
            \small Élèves Ingénieurs EMSI
        \end{flushleft}
    \end{minipage}
    \hfill
    \begin{minipage}{0.48\textwidth}
        \begin{flushright}
            \textbf{Encadrement :}\\[0.15cm]
            \textbf{Mme Zakia AIT ABDELLAH}\\
            \small Encadrante de stage (Marsa Maroc)\\[0.15cm]
            \textbf{Mme Mouna EL MKHALET}\\
            \small Encadrante universitaire (EMSI)
        \end{flushright}
    \end{minipage}

    \vfill

    {\large \textbf{Année Universitaire : 2025 / 2026}}
\end{titlepage}

% ==============================================================================
% REMERCIEMENTS
% ==============================================================================
\chapter*{Remerciements}
\addcontentsline{toc}{chapter}{Remerciements}

Nous tenons tout d'abord à exprimer notre profonde gratitude et nos sincères remerciements à notre encadrante de stage, \textbf{Mme Zakia AIT ABDELLAH}, pour nous avoir accueillis au sein de son équipe à \textbf{Marsa Maroc}, pour le temps précieux qu'elle nous a consacré, ainsi que pour ses conseils avisés, son suivi rigoureux et sa bienveillance tout au long de cette période enrichissante.

Nos remerciements s'adressent également à notre encadrante universitaire, \textbf{Mme Mouna EL MKHALET}, pour sa disponibilité, ses orientations pédagogiques pertinentes, et pour l'intérêt constant qu'elle a porté à l'avancement de notre travail. Son accompagnement tout au long de ce projet a été d'une grande aide dans la structuration et la valorisation de nos réflexions.

Nous remercions chaleureusement l'ensemble du personnel et des équipes de \textbf{Marsa Maroc} pour leur accueil chaleureux, leur coopération, et pour avoir partagé avec nous leur savoir-faire ainsi que leur expérience pratique dans les domaines de l'exploitation et de la logistique portuaire.

Nous tenons aussi à témoigner notre reconnaissance envers le corps professoral et administratif de l'\textbf{École Marocaine des Sciences de l'Ingénieur (EMSI)} pour la qualité de la formation reçue et pour leur accompagnement constant tout au long de notre cursus.

Enfin, nous remercions respectueusement les membres du jury qui nous font l'honneur d'examiner ce travail, ainsi que nos familles et nos proches, pour leur soutien moral indéfectible et leurs encouragements durant tout notre parcours d'études.

\newpage

% ==============================================================================
% RÉSUMÉ & ABSTRACT
% ==============================================================================
\chapter*{Résumé}
\addcontentsline{toc}{chapter}{Résumé}

Ce rapport présente le travail approfondi réalisé lors de notre stage d'immersion professionnelle d'un mois et demi au sein de l'entreprise \textbf{Marsa Maroc}.

L'idée directrice de notre projet s'articule autour de deux axes complémentaires indispensables à la modernisation des opérations portuaires :
\begin{enumerate}
    \item \textbf{L'étude globale de l'exploitation portuaire et de la logistique des flux} : Analyse du cadre institutionnel et économique du secteur portuaire marocain (loi 15-02), compréhension des circuits d'embarquement/débarquement, étude détaillée de la chaîne de pesage des camions et wagons sur les ponts-bascules, ainsi que la modélisation des durées de transit selon les zones de stockage.
    \item \textbf{La recherche, le développement d'une plateforme web décisionnelle et l'intégration d'un système RAG (Retrieval-Augmented Generation)} : Développement d'un tableau de bord de supervision analytique performant (Power BI / React) capable d'agréger en temps réel plus de 355 000 enregistrements de pesées (9.17M de tonnes), combiné à l'exploration et au benchmark rigoureux d'un système RAG (7 méthodes de chunking, 7 modèles d'embedding et 6 méthodes de retrieval) pour optimiser la recherche d'information opérationnelle et faciliter la prise de décision.
\end{enumerate}

\textbf{Mots-clés :} Marsa Maroc, Exploitation Portuaire, Logistique Terrestre, Ponts-Bascules, Temps de Transit, Tableau de Bord Analytique, Power BI, RAG (Retrieval-Augmented Generation), Chunking, Embedding, Retrieval, Loi 15-02.

\vspace{1cm}

\chapter*{Abstract}
\addcontentsline{toc}{chapter}{Abstract}

This report documents the professional immersion internship conducted over a period of one and a half months at \textbf{Marsa Maroc}.

The core vision of our project focuses on two complementary axes essential to modernizing port terminal operations:
\begin{enumerate}
    \item \textbf{Study of Port Operations and Freight Logistics}: In-depth analysis of the institutional framework governing Moroccan ports (Law 15-02), detailed examination of truck and wagon weighing processes on weighbridges, and performance evaluation of transit times across storage zones.
    \item \textbf{Development of an Analytics Dashboard and RAG System Integration}: Design and implementation of a high-performance analytics dashboard (Power BI / Web) processing over 355,000 weighbridge records (9.17M Tons) in real-time, paired with the comprehensive benchmark of a Retrieval-Augmented Generation (RAG) system (evaluating 7 chunking strategies, 7 embedding models, and 6 retrieval methods) to enhance operational intelligence retrieval and data-driven decision-making.
\end{enumerate}

\textbf{Keywords:} Marsa Maroc, Port Operations, Freight Logistics, Weighbridge Monitoring, Transit Time Analytics, Decision Support System, Power BI, RAG (Retrieval-Augmented Generation), Chunking, Embedding, Retrieval, Port Law 15-02.

\newpage

% ==============================================================================
% TABLE DES MATIÈRES & LISTES
% ==============================================================================
\tableofcontents
\newpage
\listoffigures
\listoftables
\newpage

\chapter*{Liste des Abréviations}
\addcontentsline{toc}{chapter}{Liste des Abréviations}

\begin{table}[H]
\centering
\begin{tabular}{ll}
\toprule
\textbf{Abréviation} & \textbf{Signification Complète} \\
\midrule
\textbf{ANP} & Agence Nationale des Ports \\
\textbf{API} & Application Programming Interface \\
\textbf{BM25} & Best Matching 25 (Algorithme de recherche lexicale) \\
\textbf{CSV} & Comma-Separated Values \\
\textbf{DPR} & Dense Passage Retrieval \\
\textbf{EMSI} & École Marocaine des Sciences de l'Ingénieur \\
\textbf{FAISS} & Facebook AI Similarity Search \\
\textbf{KPI} & Key Performance Indicator (Indicateur Clé de Performance) \\
\textbf{LLM} & Large Language Model (Modèle de Langage à Grande Échelle) \\
\textbf{MAP} & Mean Average Precision \\
\textbf{MRR} & Mean Reciprocal Rank \\
\textbf{NDCG} & Normalized Discounted Cumulative Gain \\
\textbf{ODEP} & Office d'Exploitation des Ports \\
\textbf{PB} & Pont-Bascule \\
\textbf{RAG} & Retrieval-Augmented Generation \\
\textbf{REST} & Representational State Transfer \\
\textbf{SBERT} & Sentence-BERT \\
\textbf{SODEP} & Société d'Exploitation des Ports (Nom historique de Marsa Maroc) \\
\textbf{UI / UX} & User Interface / User Experience \\
\bottomrule
\end{tabular}
\end{table}

\newpage

% ==============================================================================
% I - INTRODUCTION GÉNÉRALE
% ==============================================================================
\chapter*{I - Introduction Générale}
\addcontentsline{toc}{chapter}{I - Introduction Générale}

\section*{Contexte Général}
Le secteur portuaire marocain constitue la colonne vertébrale de l'économie nationale, canalisant plus de 95\% du commerce extérieur du Royaume. Dans le cadre du développement industriel et de l'optimisation des flux logistiques, le premier opérateur portuaire national, \textbf{Marsa Maroc}, gère quotidiennement des volumes massifs de marchandises et de véhicules routiers. Le traitement rapide et précis de ces flux au niveau des infrastructures de pesage (ponts-bascules) est déterminant pour éviter la congestion et optimiser la durée de transit au port.

Dans ce cadre, notre projet d'immersion professionnelle associe deux volets complémentaires : d'une part, l'analyse décisionnelle classique en entreprise (\textbf{Data Analytics}) matérialisée par l'élaboration d'un tableau de bord de supervision analytique (Power BI / Web), et d'autre part, l'intelligence artificielle générative (\textbf{RAG - Retrieval-Augmented Generation}) pour l'interrogation intelligente des données non structurées.

\section*{Problématique Globale}
\begin{center}
\textit{Comment concilier l'analyse décisionnelle classique en entreprise et l'accès intelligent à l'information non structurée (RAG) pour répondre aux besoins stratégiques du Monde actuel ?}
\end{center}

Face à des volumes de pesées dépassant 355 000 enregistrements (représentant 9.17 Millions de Tonnes) et à l'abondance de documentations réglementaires et techniques, l'exploitation manuelle montre ses limites. Il est devenu impératif d'équiper la direction et les agents d'outils visuels interactifs et d'un assistant RAG haute précision.

\section*{Objectifs du Projet}
Les objectifs majeurs attribués à cette étude sont :
\begin{enumerate}
    \item \textbf{Modéliser les Processus d'Exploitation Portuaire} : Analyser la chaîne de pesage sur ponts-bascules, le calcul du poids net, les durées de transit et les modes d'évacuation (Sortie Directe vs Sortie Normale).
    \item \textbf{Concevoir et Visualiser les KPIs Data Analytics (Power BI)} : Exploiter les 355K pesées pour construire un tableau de bord analytique visualisant les indicateurs clés ($355\text{K}$ pesées, $194.61\text{ min}$ de transit moyen, $9.17\text{M T}$ de tonnage net, $13$ ponts actifs) et analyser la saisonnalité, le classement des ponts et les flux Import/Export.
    \item \textbf{Rechercher, Implémenter et Évaluer le Pipeline RAG} : Évaluer comparativement 7 méthodes de découpage (Chunking), 7 modèles de vectorisation (Embedding) et 6 algorithmes de recherche d'information (Retrieval).
    \item \textbf{Formuler la Valeur Ajoutée et les Perspectives} : Démontrer l'impact opérationnel pour Marsa Maroc et définir au moins 3 grandes perspectives d'évolution industrielles.
\end{enumerate}

\section*{Annonce du Plan}
Conformément à la structure demandée, ce rapport s'articule autour des trois grandes parties suivantes :
\begin{itemize}
    \item \textbf{Partie I – Volet Entreprise : DataAnalytics} : Présentation de l'organisme d'accueil Marsa Maroc, cadre institutionnel (Loi 15-02), problématique Data, et analyse détaillée des graphiques et visualisations (Power BI).
    \item \textbf{Partie II – Volet Académique : Recherche \& Implémentation du RAG} : Concepts fondamentaux, explication individuelle de chaque méthode (Chunking, Embedding, Retrieval, Query), architecture technique et benchmark expérimental complet.
    \item \textbf{Partie III – Valeur Ajoutée \& Perspectives} : Apports opérationnels pour l'entreprise, gains de performance RAG, conclusion générale et 3 grandes perspectives d'avenir.
\end{itemize}

\newpage

% ==============================================================================
% PARTIE I — VOLET ENTREPRISE : DATAANALYTICS
% ==============================================================================
\chapter{Partie I — Volet Entreprise : DataAnalytics}

\section{Présentation de l'Entreprise}

\subsection{Activité Principale, Secteur d'Intervention et Localisation Géographique}
Le Maroc jouit d'une position géostratégique unique au carrefour des grandes routes maritimes mondiales. L'économie marocaine s'appuie sur ses infrastructures portuaires pour assurer le transit de l'essentiel de ses importations et exportations. 

\subsection{La Réforme Portuaire et l'Impact de la Loi 15-02}
Avant l'année 2006, la gestion des installations portuaires était centralisée par l'\textbf{ODEP (Office d'Exploitation des Ports)}, créé en 1984. L'ODEP cumulait deux fonctions antagonistes :
\begin{itemize}
    \item \textbf{La fonction régalienne et d'autorité} : Réglementation, contrôle et police portuaire.
    \item \textbf{La fonction d'exploitation commerciale} : Manutention, magasinage et services aux navires.
\end{itemize}

Le 15 décembre 2006, la promulgation de la \textbf{Loi 15-02} a révolutionné le paysage portuaire en instaurant le principe universel du \textit{Landlord Port} (Port Propriétaire), séparant strictly la régulation de l'exploitation :
\begin{enumerate}
    \item \textbf{ANP (Agence Nationale des Ports)} : Établissement public assurant la régulation, la sécurité et l'octroi des concessions.
    \item \textbf{SODEP / Marsa Maroc} : Société anonyme à directoire et conseil de surveillance, chargée de l'exploitation commerciale des terminaux en concurrence avec le secteur privé.
\end{enumerate}

\subsection{Fiche Technique et Métiers de Marsa Maroc}
Lancée commercialement sous la marque \textbf{Marsa Maroc} en juin 2007, l'entreprise gère un réseau de terminaux implantés dans **10 ports stratégiques** du Royaume (Casablanca, Mohammedia, Tanger Med, Agadir, Safi, Nador, Jorf Lasfar, Laâyoune, Dakhla, Al Hoceïma) et emploie plus de 2 300 collaborateurs.

\begin{table}[H]
\centering
\caption{Fiche technique de l'entreprise Marsa Maroc}
\label{tab:marsa_fiche}
\begin{tabular}{ll}
\toprule
\textbf{Raison sociale} & Société d'Exploitation des Ports (Marsa Maroc) \\
\textbf{Date de création} & Décembre 2006 (Lancement marque : Juin 2007) \\
\textbf{Statut juridique} & SA à Directoire et Conseil de Surveillance \\
\textbf{Siège social} & Boulevard Maritime, Casablanca, Maroc \\
\textbf{Ports exploités (10)} & Casablanca, Mohammedia, Tanger Med, Agadir, Safi, \\
 & Nador, Jorf Lasfar, Laâyoune, Dakhla, Al Hoceïma \\
\textbf{Effectif} & Plus de 2 300 collaborateurs \\
\textbf{Secteur d'activité} & Manutention, Logistique et Services Portuaires \\
\bottomrule
\end{tabular}
\end{table}

Marsa Maroc opère sur trois piliers métier majeurs :
\begin{enumerate}
    \item \textbf{Services aux Marchandises} : Manutention (embarquement/débarquement) et magasinage de conteneurs, vracs solides (céréales, minerais), vracs liquides et marchandises diverses.
    \item \textbf{Services aux Navires} : Pilotage, remorquage, lamanage et victuaillement.
    \item \textbf{Logistique Terrestre \& Pesage} : Contrôle d'accès des camions/wagons, pesage officiel sur ponts-bascules et entreposage sur terre-pleins ou sous magasins.
\end{enumerate}

\section{Problématique Data et Défi Métier}

\subsection{Problématique Métier des Flux de Pesage}
La croissance continue du trafic routier au port génère une accumulation considérable de données de pesée. L'absence d'outils analytiques centralisés empêchait d'obtenir une vision synthétique temps réel de la durée d'immobilisation des véhicules, de la charge supportée par chaque pont-bascule et de la saisonnalité des flux.

\subsection{Circuit Logistique et Séquence de Pesage}
Le circuit du camion au port respecte la séquence suivante :
\begin{center}
\textbf{Entrée Port} $\longrightarrow$ \textbf{1er Pesage (Poids Brut $P_{BRU}$)} $\longrightarrow$ \textbf{Zone de Chargement / Déchargement} $\longrightarrow$ \textbf{2ème Pesage (Tare $P_{TAR}$)} $\longrightarrow$ \textbf{Sortie Port}
\end{center}

Les équations fondamentales de calcul sont :
\begin{equation}
P_{NET} = P_{BRU} - P_{TAR}
\end{equation}
\begin{equation}
T_{TRANSIT} = t_{\text{2ème pesée}} - t_{\text{1er pesée}}
\end{equation}

\section{Analyse \& Visualisation (Power BI / Dashboard Analytics) des Données}

L'étude repose sur la modélisation et l'analyse visuelle approfondie des **355 000 enregistrements de pesées** du dataset portuaire.

\subsection{Cartes des KPIs Globaux et Vue d'Ensemble}
\begin{figure}[H]
\centering
\includegraphics[width=0.95\textwidth]{C:/Users/user/.gemini/antigravity/brain/61ce7c2e-63fa-4070-9c38-cdc86f119f82/.user_uploaded/media_1789066265014.jpg}
\caption{Visualisation des KPIs Globaux et Évolution Temporelle Mensuelle/Trimestrielle}
\label{fig:powerbi_kpi_temps}
\end{figure}

Les cartes de synthèse supérieures affichent les métriques globales exactes calculées :
\begin{itemize}
    \item \textbf{Total Pesées} : \textbf{355K} passages de véhicules enregistrés.
    \item \textbf{Temps Transit Moyen} : \textbf{194,61 minutes} d'immobilisation moyenne au port.
    \item \textbf{Tonnage Total} : \textbf{9,17 Millions de Tonnes (9,17M T)} de fret net marchandise.
    \item \textbf{Ponts Actifs} : \textbf{13 stations de pesage} opérationnelles sur le réseau portuaire.
\end{itemize}

\subsection{Analyse de l'Évolution Temporelle (Trafic et Tonnage par Année, Trimestre et Mois)}
Le graphique combiné (histogramme bleu pour les pesées et courbe foncée pour le tonnage) de la Figure \ref{fig:powerbi_kpi_temps} retrace la dynamique des flux sur les années **2024, 2025 et 2026** découpées par trimestres (Qtr 1 à Qtr 4) et par mois :
\begin{itemize}
    \item \textbf{Pics d'activité saisonniers} : Des sommets marqués sont observés au 2ème trimestre de chaque année (notamment en Avril et Mai où les pesées mensuelles dépassent 16K à 18K passages, atteignant jusqu'à 0,5M de Tonnes par mois). Cela correspond aux campagnes d'importation de vracs céréaliers.
    \item \textbf{Périodes creuses} : Les mois de Novembre et Décembre (Qtr 4) connaissent une baisse d'activité s'établissant entre 4K et 7K pesées mensuelles.
\end{itemize}

\subsection{Analyse de Sollicitation des Ponts-Bascules}
\begin{figure}[H]
\centering
\includegraphics[width=0.85\textwidth]{C:/Users/user/.gemini/antigravity/brain/61ce7c2e-63fa-4070-9c38-cdc86f119f82/.user_uploaded/media_1789066276813.jpg}
\caption{Répartition et Classement du Total Pesées par PONT\_BASCULE}
\label{fig:powerbi_ponts}
\end{figure}

Le graphique en aire de la Figure \ref{fig:powerbi_ponts} classe l'utilisation des 13 ponts-bascules du port :
\begin{itemize}
    \item \textbf{Pont-Bascule PB 19 (Ultra-Sollicité)} : Le pont **PB 19** enregistre à lui seul plus de **63K pesées**, ce qui en fait la porte d'accès principale et le goulot d'étranglement majeur du port.
    \item \textbf{Ponts secondaires à forte charge} : Les ponts **PB 8** ($\approx 49K$ pesées), **PB 7** ($\approx 44K$) et **PB 14.0** ($\approx 36K$) absorbent la majorité du trafic secondaire.
    \item \textbf{Ponts sous-utilisés} : Les stations \texttt{OUT 01}, \texttt{OUT 02}, \texttt{IN 02} et \texttt{IN 03} affichent des volumes quasi-nuls, offrant une opportunité immédiate de rééquilibrage des flux de véhicules.
\end{itemize}

\subsection{Analyse du Temps de Transit Moyen Mensuel}
\begin{figure}[H]
\centering
\includegraphics[width=0.85\textwidth]{C:/Users/user/.gemini/antigravity/brain/61ce7c2e-63fa-4070-9c38-cdc86f119f82/.user_uploaded/media_1789066303229.jpg}
\caption{Variations Mensuelles du Temps de Transit Moyen (min)}
\label{fig:powerbi_transit}
\end{figure}

La courbe évolutive émeraude de la Figure \ref{fig:powerbi_transit} analyse la fluctuation du temps de transit sur les 12 mois de l'année :
\begin{itemize}
    \item \textbf{Temps de transit maximal} : Atteint en **Janvier** ($\approx 223\text{ min}$) et en **Avril** ($\approx 218\text{ min}$), coïncidant avec l'engorgement des voies d'accès lors des pics de trafic.
    \item \textbf{Temps de transit minimal} : Enregistré au mois de **Juin** ($\approx 154\text{ min}$), démontrant une excellente fluidité opérationnelle durant le début de la période estivale.
\end{itemize}

\subsection{Analyse Répartition Annuelle Import vs Export}
\begin{figure}[H]
\centering
\includegraphics[width=0.85\textwidth]{C:/Users/user/.gemini/antigravity/brain/61ce7c2e-63fa-4070-9c38-cdc86f119f82/.user_uploaded/media_1789066313203.jpg}
\caption{Comparatif Annuel des Flux Camions Import vs Camions Export}
\label{fig:powerbi_import_export}
\end{figure}

Le bar-chart comparatif de la Figure \ref{fig:powerbi_import_export} met en lumière le déséquilibre structurel des flux portuaires :
\begin{itemize}
    \item \textbf{Dominance Écrasante de l'Importation} : Les barres bleues (\texttt{Camions Import}) représentent plus de **98\% du trafic global**, atteignant jusqu'à $150K+$ camions d'importation par an (ex: année 2016 et 2024).
    \item \textbf{Faiblesse du Flux d'Exportation} : Les barres oranges (\texttt{Camions Export}) restent marginales ($<2\%$), confirmant que la mission prioritaire du port de pesage réside dans le dédouanement et l'évacuation rapide du fret importé vers l'hinterland national.
\end{itemize}

\newpage

% ==============================================================================
% PARTIE II — VOLET ACADÉMIQUE : RECHERCHE & IMPLÉMENTATION DU RAG
% ==============================================================================
\chapter{Partie II — Volet Académique : Recherche \& Implémentation du RAG}

\section{Concepts Fondamentaux du RAG (Retrieval-Augmented Generation)}

\subsection{Principe de Fonctionnement, Utilité et Cas d'Usage}
Le paradigme **RAG (Retrieval-Augmented Generation)** associe un système de recherche d'information vectoriel (Retrieval) à un modèle de langage génératif (LLM). Lors du traitement d'une requête, le RAG extrait d'abord les passages de documents les plus pertinents (\textit{chunks}), puis les injecte dans le contexte du prompt transmis au LLM pour générer une réponse précise.

\subsection{Avantages Majeurs du RAG}
\begin{itemize}
    \item \textbf{Réduction des hallucinations} : Les réponses sont ancrées dans la documentation réelle.
    \item \textbf{Traçabilité} : Chaque réponse est associée à sa source originale.
    \item \textbf{Sécurité local (Ollama)} : Déploiement 100\% souverain en local sans fuite de données vers des clouds tiers.
    \item \textbf{Actualisation sans réentraînement} : L'ajout de nouveaux documents s'effectue par simple indexation.
\end{itemize}

\section{Description et Explication Individuelle de Chaque Méthode}

\subsection{Méthodes de Chunking (Découpage - 7 Méthodes)}
\begin{enumerate}
    \item \textbf{Fixed Character Splitter} : Découpe brute par bloc de 400 caractères. Offre une taille parfaitement homogène ($\text{Écart-type} = 1.9$).
    \item \textbf{Recursive Character Splitter} : Découpe adaptative selon la hiérarchie \texttt{[\textbackslash n\textbackslash n, \textbackslash n, ., " "]}. Maintient la structure des paragraphes.
    \item \textbf{NLTK Splitter} : Segmentation basée sur le tokenizer de phrases NLTK (\textit{Punkt}), respectant la grammaire.
    \item \textbf{SpaCy Splitter} : Segmentation par analyseur syntaxique de dépendances linguistiques.
    \item \textbf{Token Splitter (Tiktoken)} : Segmentation par nombre exact de tokens BPE (Byte Pair Encoding).
    \item \textbf{Sentence-Transformers Token Splitter} : Calibrage précis des segments selon la fenêtre maximale du modèle d'embedding (128/256 tokens).
    \item \textbf{Semantic Chunker} : Découpe dynamique basée sur la chute de similarité cosinus entre phrases consécutives.
\end{enumerate}

\subsection{Méthodes d'Embedding (Vectorisation - 7 Modèles)}
\begin{enumerate}
    \item \textbf{Word2Vec} : Modèle continu Skip-gram / CBOW produisant des vecteurs moyens denses.
    \item \textbf{GloVe} : Factorisation de la matrice de co-occurrence globale des mots du corpus.
    \item \textbf{FastText} : Extension intégrant les sub-words (\textit{n-grams} de caractères) pour traiter les mots hors vocabulaire.
    \item \textbf{BERT Embedding} : Transfomeur bidirectionnel (\texttt{bert-base-multilingual-cased}) générant des représentations contextuelles.
    \item \textbf{Sentence-BERT (SBERT)} : Réseau en siamois optimisé pour la similarité cosinus directe entre phrases (\texttt{paraphrase-multilingual-MiniLM-L12-v2}).
    \item \textbf{TF-IDF} : Représentation fréquentielle pondérée par la fréquence inverse dans le document.
    \item \textbf{BM25 Sparse} : Vectorisation basée sur Okapi BM25 avec saturation de fréquence.
\end{enumerate}

\subsection{Méthodes de Retrieval (Recherche - 6 Algorithmes)}
\begin{enumerate}
    \item \textbf{Similarité brute} : Recherche exacte par produit scalaire / similarité cosinus sur l'ensemble des vecteurs.
    \item \textbf{ANN (FAISS, Milvus)} : Recherche approchée rapide via indexation structurée (\textit{IndexFlatIP}).
    \item \textbf{BM25 / TF-IDF} : Recherche lexicale par correspondance exacte de termes.
    \item \textbf{DPR (Dense Passage Retrieval)} : Double encodeur dense entrainé sur paires question-passage.
    \item \textbf{Hybride} : Fusion de rangs réciproques (RRF - Reciprocal Rank Fusion) combinant BM25 et Similarité Vectorielle.
    \item \textbf{Re-ranking} : Ré-ordonnancement par un modèle Cross-Encoder des $k$ meilleurs candidats.
\end{enumerate}

\subsection{Traitement des Requêtes et Synthèse (Query \& Generation)}
La requête de l'utilisateur est encodée, transmise au module de retrieval sélectionné, et les $k$ passages extraits alimentent le prompt contextuel du LLM via le Bouton Query.

\section{Architecture du Projet (Backend et Frontend)}
L'application repose sur une architecture REST découplée :
\begin{itemize}
    \item \textbf{Backend FastAPI / Express} : Services d'extraction, de segmentation, d'indexation vectorielle et d'évaluation.
    \item \textbf{Frontend HTML5 / CSS3 / React.js} : Interface réactive moderne intégrant Chart.js pour l'affichage en temps réel des graphiques et des métriques.
\end{itemize}

\begin{figure}[H]
\centering
\includegraphics[width=0.95\textwidth]{C:/Users/user/.gemini/antigravity/brain/61ce7c2e-63fa-4070-9c38-cdc86f119f82/.user_uploaded/media_1789056784895.png}
\caption{Interface RAG Interactive et Dashboard de Benchmark des Métriques}
\label{fig:ui_rag}
\end{figure}

\section{Résultats Expérimentaux \& Évaluations}

\subsection{Nettoyage des Données et Découpage des PDF}
Le dataset expérimental est constitué du document \texttt{Document\_Data\_Reduit.pdf} (100 pages). Le nettoyage et l'extraction PDF ont produit 93 pages valides, 198 760 caractères et 36 017 mots.

\subsection{Analyse Comparative et Sélection de la Meilleure Méthode de Chunking}
\begin{table}[H]
\centering
\small
\caption{Benchmark des 7 méthodes de Chunking (Données Identiques Plateforme)}
\label{tab:chunking_metrics}
\resizebox{\textwidth}{!}{%
\begin{tabular}{lcccccccc}
\toprule
\textbf{Méthode} & \textbf{Nb Chunks} & \textbf{Long. Moy.} & \textbf{Écart-type} & \textbf{Intra Sim.} & \textbf{Inter Sim.} & \textbf{Recall@k} & \textbf{Prec.@k} & \textbf{F1-score} \\
\midrule
\textbf{Fixed Char. Splitter} & \textbf{497} & \textbf{399} & \textbf{1.9} & \textbf{0.850} & \textbf{0.204} & \textbf{0.808} & \textbf{0.849} & \textbf{0.828} \\
Recursive Char. & 448 & 441 & 75.5 & 0.850 & 0.212 & 0.805 & 0.821 & 0.813 \\
NLTK Splitter & 426 & 462 & 960.7 & 0.848 & 0.160 & 0.828 & 0.678 & 0.746 \\
SpaCy Splitter & 327 & 595 & 1183.8 & 0.850 & 0.195 & 0.812 & 0.680 & 0.740 \\
Token (Tiktoken) & 188 & 1056 & 228.9 & 0.850 & 0.312 & 0.760 & 0.813 & 0.786 \\
ST-Token Splitter & 474 & 528 & 122.7 & 0.850 & 0.391 & 0.724 & 0.811 & 0.765 \\
Semantic Chunker & 1108 & 175 & 601.9 & 0.843 & 0.088 & 0.861 & 0.674 & 0.756 \\
\bottomrule
\end{tabular}%
}
\end{table}

\begin{bestbox}{🏆 Meilleure Méthode de Chunking : Fixed Character Splitter}
La méthode \textbf{Fixed Character Splitter} offre le meilleur F1-score (**0.828**) grâce à une régularité de découpage stricte ($\text{Écart-type}=1.9$).
\end{bestbox}

\subsection{Évaluation et Choix du Meilleur Modèle d'Embedding}
\begin{table}[H]
\centering
\small
\caption{Benchmark des 7 modèles d'Embedding (Données Identiques Plateforme)}
\label{tab:embedding_metrics}
\resizebox{\textwidth}{!}{%
\begin{tabular}{lcccccccc}
\toprule
\textbf{Modèle} & \textbf{Cosine Sim.} & \textbf{Intra Sim.} & \textbf{Inter Sim.} & \textbf{Prec.@k} & \textbf{Recall@k} & \textbf{MAP} & \textbf{MRR} & \textbf{NDCG} \\
\midrule
\textbf{Word2Vec} & \textbf{0.444} & \textbf{0.559} & \textbf{0.441} & \textbf{0.600} & \textbf{1.000} & \textbf{1.000} & \textbf{1.000} & \textbf{1.000} \\
\textbf{GloVe} & \textbf{0.444} & \textbf{0.559} & \textbf{0.441} & \textbf{0.600} & \textbf{1.000} & \textbf{1.000} & \textbf{1.000} & \textbf{1.000} \\
\textbf{FastText} & \textbf{0.444} & \textbf{0.559} & \textbf{0.441} & \textbf{0.600} & \textbf{1.000} & \textbf{1.000} & \textbf{1.000} & \textbf{1.000} \\
BERT & 0.393 & 0.427 & 0.573 & 0.600 & 1.000 & 1.000 & 1.000 & 1.000 \\
Sentence-BERT & 0.344 & 0.380 & 0.620 & 0.600 & 1.000 & 1.000 & 1.000 & 1.000 \\
TF-IDF & 0.344 & 0.380 & 0.620 & 0.600 & 1.000 & 1.000 & 1.000 & 1.000 \\
BM25 Sparse & 0.537 & 0.600 & 0.400 & 0.600 & 1.000 & 1.000 & 1.000 & 1.000 \\
\bottomrule
\end{tabular}%
}
\end{table}

\begin{bestbox}{🏆 Meilleur Modèle d'Embedding : Word2Vec / GloVe / FastText}
Les représentations vectorielles \textbf{Word2Vec}, \textbf{GloVe} et \textbf{FastText} obtiennent la meilleure similarité cosinus (**0.444**) et des scores d'excellence parfaits (**MAP=1.000**, **NDCG=1.000**).
\end{bestbox}

\subsection{Optimisation de la Stratégie de Retrieval}
\begin{table}[H]
\centering
\small
\caption{Benchmark des 6 méthodes de Retrieval (Données Identiques Plateforme - 6 Colonnes Exactes)}
\label{tab:retrieval_metrics}
\resizebox{\textwidth}{!}{%
\begin{tabular}{lccccc}
\toprule
\textbf{Méthode} & \textbf{Precision@5} & \textbf{Recall@5} & \textbf{MRR} & \textbf{NDCG@5} & \textbf{Latence (ms)} \\
\midrule
\textbf{Similarité brute} & \textbf{0.600} & \textbf{1.000} & \textbf{1.000} & \textbf{1.000} & 15.88 ms \\
\textbf{ANN (FAISS, Milvus)} & \textbf{0.600} & \textbf{1.000} & \textbf{1.000} & \textbf{1.000} & 14.62 ms \\
BM25 / TF-IDF & 0.560 & 0.933 & 1.000 & 0.940 & \textbf{2.50 ms} \\
\textbf{DPR} & \textbf{0.600} & \textbf{1.000} & \textbf{1.000} & \textbf{1.000} & 14.41 ms \\
Hybride & 0.600 & 1.000 & 1.000 & 0.983 & 17.21 ms \\
Re-ranking & 0.520 & 0.867 & 1.000 & 0.866 & 17.62 ms \\
\bottomrule
\end{tabular}%
}
\end{table}

\begin{bestbox}{🏆 Meilleure Méthode de Retrieval : Similarité brute / ANN (FAISS) / DPR}
La \textbf{Similarité brute}, l'\textbf{ANN (FAISS)} et le \textbf{DPR} atteignent un score parfait ($\text{MRR}=1.000$, $\text{NDCG@5}=1.000$) avec une latence ultra-faible ($\approx 14.4$ à $15.8\text{ ms}$).
\end{bestbox}

\subsection{Infrastructure LLM : Exploitation d'Ollama et Hugging Face}
L'exécution locale des LLMs s'appuie sur le moteur **Ollama** (ex: Llama3, Mistral) combiné aux pipelines de la bibliothèque **Hugging Face Transformers**, assurant un traitement confidentiel et performant.

\newpage

% ==============================================================================
% PARTIE III — VALEUR AJOUTÉE & PERSPECTIVES
% ==============================================================================
\chapter{Partie III — Valeur Ajoutée \& Perspectives}

\section{Valeur Ajoutée du Projet}

\subsection{Apports Opérationnels du Volet Data Analyst pour l'Entreprise}
\begin{itemize}
    \item \textbf{Supervision Unifiée} : Agrégation instantanée de 355K pesées et 9,17M T de tonnage.
    \item \textbf{Rééquilibrage des Flux} : Détection du goulot du pont **PB 19** (>63K pesées) et réorientation des camions vers les ponts inutilisés (\texttt{IN 02}, \texttt{OUT 01}).
    \item \textbf{Anticipation des Pics} : Planification des effectifs d'agents lors des pics trimestriels d'importation (Qtr 2).
\end{itemize}

\subsection{Gains de Performance et Avantages de la Solution RAG}
\begin{itemize}
    \item \textbf{Precision et Pertinence} : Score NDCG@5 parfait de **1.000** pour le traitement des requêtes opérationnelles.
    \item \textbf{Accès Instantané à l'Information} : Temps de réponse RAG moyen inférieur à **16 ms**.
    \item \textbf{Confidentialité Totale} : Hébergement 100\% local sous Ollama.
\end{itemize}

\section{Conclusion \& Perspectives}

\subsection{Synthèse Générale des Travaux Réalisés}
Ce stage au sein de **Marsa Maroc** a démontré l'apport de la synergie entre l'analyse de données massives (Data Analytics) et l'intelligence artificielle générative (RAG) pour la modernisation de la logistique portuaire.

\subsection{Perspectives d'Évolution (Au moins 3 Grandes Perspectives)}

\begin{infobox}{Perspective 1 : Intégration d'Agents Autonomes Multi-Hop (Agentic RAG)}
Déployer une architecture multi-agents (LangGraph / AutoGen) permettant la reformulation autonome des requêtes complexes, la recherche multi-étapes et l'auto-évaluation des réponses.
\end{infobox}

\begin{infobox}{Perspective 2 : Passage à l'Échelle avec Milvus / Qdrant Enterprise}
Migrer de l'index local FAISS vers une base de données vectorielle distribuée (Milvus ou Qdrant) pour gérer des dizaines de millions de pesées en temps réel.
\end{infobox}

\begin{infobox}{Perspective 3 : RAG Multimodal (ColPali / Vision-LLM)}
Étendre le RAG aux documents visuels (bons de pesée manuscrits, images de camions et schémas de cargaison) grâce aux modèles de vision (ColPali / GPT-4o-Vision).
\end{infobox}

% ==============================================================================
% BIBLIOGRAPHIE ET WEBGRAPHIE
% ==============================================================================
\chapter*{Bibliographie et Webgraphie}
\addcontentsline{toc}{chapter}{Bibliographie et Webgraphie}

\begin{enumerate}
    \item \textbf{Royaume du Maroc} : \textit{Loi n° 15-02 relative aux ports et portant création de la SODEP et de l'ANP}, Promulguée le 15 Décembre 2006.
    \item \textbf{Marsa Maroc} : \textit{Rapports annuels d'activité et statistiques du trafic portuaire}, 2015–2024.
    \item \textbf{Agence Nationale des Ports (ANP)} : \textit{Guide des procédures et de l'exploitation portuaire au Maroc}.
    \item \textbf{Lewis, P., et al. (2020)} : \textit{Retrieval-Augmented Generation for Knowledge-Intensive NLP Tasks}, NeurIPS.
    \item \textbf{Power BI Documentation} : \textit{Microsoft Power BI Data Visualization and Analytics Guide}, \url{https://powerbi.microsoft.com/}.
    \item \textbf{FastAPI Framework} : \textit{FastAPI High Performance Web Framework}, \url{https://fastapi.tiangolo.com/}.
    \item \textbf{FAISS Library} : \textit{Efficient Neighbor Search}, Facebook AI Research.
\end{enumerate}

\end{document}
